% The cube as a 2-lift of K_4: each fibre is a pair of antipodal vertices.
% Requires: \liftunit (unit length), colours ink and accent.
\begin{tikzpicture}[x=\liftunit, y=\liftunit,
    e/.style={black!55, semithick},
    f1/.style={circle, draw=ink, fill=ink, text=white, inner sep=1.4pt, font=\footnotesize},
    f2/.style={circle, draw=accent, fill=accent, text=white, inner sep=1.4pt, font=\footnotesize},
    f3/.style={circle, draw=ink, fill=white, text=ink, inner sep=1.4pt, font=\footnotesize},
    f4/.style={circle, draw=ink, fill=black!18, text=ink, inner sep=1.4pt, font=\footnotesize}]
  % base graph K_4
  \begin{scope}
    \node[f1] (a) at (0,0) {1}; \node[f2] (b) at (2,0) {2};
    \node[f4] (c) at (2,2) {4}; \node[f3] (d) at (0,2) {3};
    \draw[e] (a)--(b)--(c)--(d)--(a) (a)--(c) (b)--(d);
    \node[font=\footnotesize, anchor=north] at (1,-0.5) {base graph $K_4$};
  \end{scope}
  \draw[->, black!55] (3.0,1) -- (4.4,1);
  % the cube: vertices are binary triples, antipodal pairs form the fibres
  \begin{scope}[xshift={5.4*\liftunit}, yshift={-0.6*\liftunit}]
    \node[f1] (v000) at (0,0) {1};     \node[f2] (v100) at (3.2,0) {2};
    \node[f4] (v110) at (3.2,3.2) {4}; \node[f3] (v010) at (0,3.2) {3};
    \node[f4] (v001) at (1,1) {4};     \node[f3] (v101) at (2.2,1) {3};
    \node[f1] (v111) at (2.2,2.2) {1}; \node[f2] (v011) at (1,2.2) {2};
    \draw[e] (v000)--(v100)--(v110)--(v010)--(v000);
    \draw[e] (v001)--(v101)--(v111)--(v011)--(v001);
    \draw[e] (v000)--(v001) (v100)--(v101) (v110)--(v111) (v010)--(v011);
    \node[font=\footnotesize, anchor=north] at (1.6,-0.5) {a $2$-lift: the cube};
  \end{scope}
\end{tikzpicture}
